%%%%%%%%%%%%%%%%%%%%%%%%%%%%%%%%%%%%%%%%%%%%%%%%%%%%%%%%%%%%%%%%%%%%%%%%%%%%%%%%
% ScaleNav -- ICRA-style manuscript (revised draft)
% Compile with ./build.sh
% Figures under pics/ are rendered when present; otherwise a labeled
% placeholder box is shown so the manuscript compiles standalone.
%%%%%%%%%%%%%%%%%%%%%%%%%%%%%%%%%%%%%%%%%%%%%%%%%%%%%%%%%%%%%%%%%%%%%%%%%%%%%%%%

\documentclass[letterpaper,10pt,conference]{ieeeconf}
\IEEEoverridecommandlockouts
\overrideIEEEmargins

\usepackage{amsmath}
\usepackage{amssymb}
\usepackage{booktabs}
\usepackage{graphicx}
\usepackage{xcolor}
\usepackage{placeins}
\usepackage{xstring}

\newcommand{\vect}[1]{\boldsymbol{#1}}
\newcommand{\mat}[1]{\mathbf{#1}}
\newcommand{\R}{\mathbb{R}}
\newcommand{\blank}{--}

% Render the figure when available; otherwise insert a labeled placeholder so
% that the manuscript compiles without the pics/ directory.
\newcommand{\safefig}[2][]{%
  \IfFileExists{#2}{\includegraphics[#1]{#2}}%
  {\begingroup\setlength{\fboxsep}{4pt}%
   \framebox[\dimexpr\linewidth-2\fboxsep-2\fboxrule]{%
     \parbox{0.88\linewidth}{\centering\vspace{1.6em}%
       \footnotesize\textit{[Figure placeholder]}\\[2pt]%
       \begingroup\footnotesize\ttfamily
       \edef\safefigpath{\detokenize{#2}}%
       \StrSubstitute{\safefigpath}{/}{/\allowbreak}[\safefigpathA]%
       \StrSubstitute{\safefigpathA}{_}{\_\allowbreak}[\safefigpathB]%
       \safefigpathB\endgroup%
       \vspace{1.6em}}}\endgroup}}

\title{\LARGE \bf
ScaleNav: Semantic Foresight on Persistent Free-Space Graphs\\
for Language-Guided Aerial Navigation
}

\author{First Author$^{1}$ and Second Author$^{2}$%
\thanks{This work was not supported by any organization.}%
\thanks{$^{1}$First Author is with the Affiliation,
        {\tt\small first.author@example.com}}%
\thanks{$^{2}$Second Author is with the Affiliation,
        {\tt\small second.author@example.com}}%
}

\IEEEaftertitletext{%
  \begin{center}
  \safefig[width=0.98\textwidth]{pics/teaser.pdf}
  \vspace{3pt}
  \parbox{0.98\textwidth}{\footnotesize
  \textbf{Fig. 1.}~ScaleNav in Map2 at synchronized evidence time $t^*$.
  (a)~RGB, semantic score, and depth show semantic response where
  depth is clipped beyond $20$~m. (b)~A shallow aerial view of the mission over
  surface-sampled UE mesh truth. Graph-node fill encodes each node's semantic
  score; crosses mark far-field semantic-node
  positions that lie within the planner's three-dimensional influence radius
  of this graph slice. The persistent graph is shown in a $\pm2$~m
  navigation-height slice. The
  graph/path state is within $90$~ms of $t^*$.
  \texttt{mission\_goal}, \texttt{frontier\_goal}, and \texttt{local\_goal} denote the task,
  route, and execution layers. Semantic risk changes verified-edge costs;
  crosses denote provisional frontiers beyond the measured backbone.}
  \end{center}%
  \setcounter{figure}{1}%
}

\begin{document}
\maketitle
\thispagestyle{empty}
\pagestyle{empty}

%%%%%%%%%%%%%%%%%%%%%%%%%%%%%%%%%%%%%%%%%%%%%%%%%%%%%%%%%%%%%%%%%%%%%%%%%%%%%%%%
\begin{abstract}
Vision-based local planners react quickly to observed geometry, but their lack
of persistent state prevents them from preserving long-horizon route
decisions; open-vocabulary perception can expose semantically meaningful
hazards well beyond reliable depth, yet it runs too slowly to shape control
directly.  This paper presents ScaleNav, a layered aerial navigation
framework that reconciles these timescales by explicitly separating reusable
mapping, language-conditioned routing, and local execution.  A
policy-independent mapper maintains a bounded, persistent graph of verified
free-space bubbles linked by collision-checked witness paths, built from depth
and odometry alone.  A slower route layer projects open-vocabulary image
responses into provisional far-field frontiers, scores continuous semantic
risk along complete witness polylines, and preserves the accepted route in
world coordinates across asynchronous graph replacement.  The selected witness
reaches the local executor through the executor's own native interface---a
frontier goal and a tracked local goal---so the demonstrated YOPO policy runs
with its released weights and architecture, without any modification or
retraining.  Throughout the stack, semantic
evidence reshapes route preference but never certifies free space: geometric
verification and depth-based local execution remain authoritative.  In 140-m
closed-loop missions, ScaleNav completes all ten trials at
$4.439\pm0.118$~m/s---$18.1\%$ faster in mean flight time than the same
unmodified executor driven by a reactive frontier goal at comparable path
length and with no loss of its perfect safety record, isolating the gain to
the route layer alone.
\end{abstract}

%%%%%%%%%%%%%%%%%%%%%%%%%%%%%%%%%%%%%%%%%%%%%%%%%%%%%%%%%%%%%%%%%%%%%%%%%%%%%%%%
\section{INTRODUCTION}
Flying quickly through an unknown environment forces a robot to reconcile
processes that naturally operate at very different rates.  A local policy must
convert the current depth image and vehicle state into a dynamically feasible
command within milliseconds
\cite{rappids,nanomap,egoplanner,yopo,loquercio}, whereas building a reusable
map, reasoning over long-horizon alternatives, and running open-vocabulary
perception all take orders of magnitude longer
\cite{octomap,voxblox,clip}.  Learned one-stage planners
\cite{yopo,loquercio} and modern control stacks \cite{kaufmann2023} have shown
how much a fast reactive layer alone can achieve; yet a single observation
cannot preserve an early left--right commitment around a long obstacle, and
once the fork disappears from view a memoryless policy is free to undo it.
Topological memories in ground navigation have repeatedly demonstrated that
retaining route structure, rather than reconstructing it from scratch at every
tick, is what converts local reactivity into reliable long-range behavior
\cite{sptm,ving,viking,vint}.

A complementary opportunity has emerged on the perception side.  RGB carries
evidence that depth does not: a power line, a glass facade, or a distant
building can be semantically obvious in pixels while remaining invisible to a
depth channel clipped at tens of meters.  Open-vocabulary models make this
evidence addressable in free language
\cite{clip,lseg,openseg,pearl}, and recent navigation systems already exploit
it to choose \emph{where to go} \cite{vlfm,lmnav,semexp,rayfronts,wildos}.
What they largely do not address is \emph{where not to go} at flight speed:
their semantic layers run at a few hertz, hold no executable notion of free
space, and hand a reactive executor a goal rather than a route.  For an aerial
vehicle, the resulting gap between a slow semantic opinion and a fast
commitment is precisely where missions are lost.

ScaleNav closes this gap by making the hierarchy explicit rather than
implicit in a monolithic network (Fig.~\ref{fig:architecture}).  Its first
layer is a plug-and-play persistent mapper: generic depth and odometry are
converted into a bounded graph of verified free-space bubbles connected by
collision-checked witness polylines, with no occupancy grid, no ESDF, and no
dependence on the downstream trajectory policy
\cite{bubble,epic,sfc}.  Its second, slower layer performs language-conditioned
route planning on that graph.  Verified bubbles form the geometric backbone;
depth-clipped semantic observations are admitted as provisional far-field
frontiers and as a continuous risk field measured along entire witnesses, not
only at endpoints.  The accepted route is stored as a world-frame witness and
remapped after each atomic graph replacement, so route identity survives the
mapper's update cycle.  Its third, fast layer is deliberately not ours: the
route layer drives the local executor only through the executor's own native
interface---a frontier goal and a tracked local goal on the accepted
witness---so the demonstrated YOPO policy \cite{yopo} runs unmodified, with
its released weights, architecture, and collision behavior, and retains full
responsibility for executable motion.

This separation lets us state the safety boundary precisely, and we regard
that boundary as a first-class design contribution.  The mapping layer
validates measured free space.  A provisional semantic-frontier connection is
an optimistic route hypothesis: it is checked against all currently known
obstacles before publication, but it is never treated as proof that
unobserved space is free.  Semantic risk repels or attracts route choices,
while current depth and the local executor remain solely responsible for
newly observed obstacles.  In the language of risk-aware planning
\cite{majumdar,blackmore}, semantic evidence enters as a shaped preference
over verified alternatives, not as a probabilistic safety guarantee.

The main contributions are:
\begin{itemize}
  \setlength\itemsep{0pt}
  \item \textbf{A plug-and-play persistent mapping layer.}
  A bounded free-space graph turns generic depth and odometry into reusable,
  witness-verified route structure without an occupancy grid or ESDF, and
  without any assumption about the local policy architecture
  (Sec.~\ref{sec:topology}).
  \item \textbf{A slow semantic route layer.}
  Provisional far-field frontiers and continuous witness-level risk give the
  planner early, language-conditioned route preference, while a world-frame
  witness with a hysteretic gate preserves route identity across asynchronous
  graph updates (Sec.~\ref{sec:semantics}).
  \item \textbf{Executor-agnostic route transfer.}
  The accepted route reaches the local executor only through the executor's
  native goal interface: no network modification, no retraining, no new loss
  terms.  We demonstrate the contract on the released YOPO policy and
  demonstrate that the gain isolates to the route layer
  (Sec.~\ref{sec:execution}).
  \item \textbf{A planner-agnostic evaluation of beyond-range semantic
  mapping.}  Building on the insight of \cite{rayfronts} that a mapping
  layer's utility can be measured without a planner in the loop, we evaluate
  our far-field semantic memory with a route-agnostic protocol---route-change
  rate, query compliance, decision distance, switch count, and semantic
  exposure under frozen geometric snapshots (Sec.~\ref{sec:experiments})---
  separating the value of the semantic representation from the behavior of
  any single planner, exactly as \cite{rayfronts}'s search-volume metric does
  for exploration.
\end{itemize}

The evaluation deliberately separates system integration, local route
adherence, and graph persistence, because each layer makes a different claim
and deserves a different falsification.  All reported semantic observations
use PEARL \cite{pearl}; the interface itself is model-agnostic.

\begin{figure*}[t]
\centering
\safefig[width=0.98\textwidth]{pics/system_architecture_paper.pdf}
\caption{Layered ScaleNav architecture.  The plug-and-play mapping layer
converts depth and odometry into persistent verified bubble topology.  The
slow route layer combines this backbone with text-conditioned risk and
provisional semantic frontiers, then preserves an accepted world-frame
witness.  The fast execution layer receives a frontier goal and an ordered
bubble corridor.  Solid geometric witnesses are verified by the mapper; a
provisional semantic connection is an optimistic frontier hypothesis, not a
free-space certificate.}
\label{fig:architecture}
\end{figure*}

%%%%%%%%%%%%%%%%%%%%%%%%%%%%%%%%%%%%%%%%%%%%%%%%%%%%%%%%%%%%%%%%%%%%%%%%%%%%%%%%
\section{RELATED WORK}
\subsection{High-Speed Local Flight}
Reactive and local planners trade global consistency for latency.  RAPPIDS
queries a raw depth image directly \cite{rappids}, NanoMap retains a short
range history for uncertainty-aware proximity queries \cite{nanomap},
EGO-Planner removes the ESDF from the optimization loop \cite{egoplanner},
and learned policies map observations to motion primitives in a single stage
\cite{yopo,loquercio,kaufmann2023}.  A parallel optimization line combines
online search with trajectory refinement: Fast-Planner and FASTER interleave
kinodynamic search with B-spline optimization \cite{fastplanner,faster}, and
SUPER recently demonstrated safety-assured point-cloud corridor flight above
$20$~m/s \cite{super}.  Free-space corridor representations underpin much of
this progress, from polyhedral safe flight corridors \cite{sfc} and receding
bubble corridors \cite{bubble} to whole-trajectory optimization within
corridor constraints \cite{gcopter}.  ScaleNav builds on this lineage but does
not compete with it: we neither claim nor reproduce SUPER's safety assurance
or flight speed.  Instead, we add the orthogonal ingredient these systems
lack---a persistent, language-conditioned route memory above the local
executor---and we keep the executor's primitive machinery unchanged, so that
a stronger local layer can be substituted beneath the route layer through an
adapter or retraining rather than a redesign
\cite{super,egoplanner,gcopter}; the substitution cost is made explicit in
Sec.~\ref{sec:discussion}.

\subsection{Persistent and Topological Spatial Memory}
The value of spatial memory for navigation predates learned policies.
Frontier-based exploration established incremental free-space reasoning
\cite{yamauchi}, extended to hierarchical roadmaps and incremental frontier
structures \cite{fuel,tare,epic}.  Metric memories store probabilistic
occupancy \cite{octomap} or distance fields \cite{voxblox}; metric-semantic
systems attach object- and scene-level meaning to such maps
\cite{kimera,hydra}.  In the learning literature, semi-parametric topological
memory \cite{sptm} and its successors \cite{ving,viking,vint,nomad} showed
that a graph of past observations supports kilometer-scale visual navigation
when paired with a capable local controller, and hierarchical compositions of
local control with topological graphs \cite{meng2020} formalize the same
division we adopt.  What these memories share is that their stored costs
remain geometric or appearance-based.  ScaleNav contributes a complementary
memory: a selected \emph{witness route} is preserved in world coordinates
across graph replacement, open-vocabulary risk is attached to verified
witnesses, and beyond-range observations are represented explicitly as
provisional frontiers rather than being silently dropped or hallucinated as
free space.

\subsection{Open-Vocabulary Scene Representation}
Vision--language pretraining made image regions addressable by free text
\cite{clip}, and dense variants bring this alignment to pixels
\cite{lseg,openseg,pearl}.  These responses have been lifted into 3-D through
feature distillation into implicit or explicit scene representations
\cite{lerf,clipfields,conceptfusion,openscene}, into queryable value maps
\cite{vlmaps}, and into object-centric or hierarchical scene graphs
\cite{conceptgraphs,hovsg}.  Most of these maps target retrieval, grounding,
or task planning on ground robots, where a wrong semantic association costs a
detour.  Aerial flight is less forgiving: the semantic layer must respect a
depth channel that is simultaneously the authoritative safety sensor and
blind beyond tens of meters.  ScaleNav therefore keeps the semantic model
behind a narrow, replaceable contract---a calibrated patch heatmap and a
confidence value---and demonstrates the contract with PEARL \cite{pearl},
while delegating all free-space certification to geometry.

\subsection{Language-Conditioned Navigation and Exploration}
Language has been used to bias exploration and goal selection at increasing
levels of abstraction: persistent semantic spatial memories for instruction
following \cite{blukis}, goal-oriented semantic exploration \cite{semexp},
frontier selection from vision--language value maps \cite{vlfm}, compositional
pipelines of language, vision, and action models \cite{lmnav}, and beyond-range
semantic ray frontiers that retain observations past the depth limit
\cite{rayfronts}.  In aerial settings, vision-and-language navigation
benchmarks \cite{aerialvln} and open-vocabulary search systems
\cite{wildos} target finding described targets, and semantic cues have guided
active perception for UAVs \cite{bartolomei}.  ScaleNav addresses the dual
problem: \emph{avoidance}.  Risk must repel a route continuously along its
witness---including before depth reaches the observation---and must do so
without ever masquerading as a geometric certificate, a distinction that
goal-seeking systems do not need to enforce.

\subsection{Risk-Aware Route Choice}
Attaching cost to uncertainty has a long history in planning under
uncertainty, from chance-constrained programs \cite{blackmore} to axiomatic
risk metrics \cite{majumdar} and risk-shaped safety corridors \cite{rast}.
We borrow the shaping principle rather than the probabilistic machinery:
semantic confidence is a monotone preference signal, not a calibrated
collision probability, so ScaleNav applies a bounded barrier that grows
sharply near unit risk while leaving geometric validity to the witness checks
(Sec.~\ref{sec:semantics}).  This keeps the semantic layer honest about what
it knows, and keeps the safety case entirely on the measured geometry.

%%%%%%%%%%%%%%%%%%%%%%%%%%%%%%%%%%%%%%%%%%%%%%%%%%%%%%%%%%%%%%%%%%%%%%%%%%%%%%%%
\section{LAYERED SYSTEM OVERVIEW}
Let the vehicle state at time $t$ be
$\vect{x}_t=(\vect{p}_t,\vect{v}_t,\vect{a}_t,\vect{R}_t)$.
The sensors provide an RGB image $I_t$, a depth image $D_t$, and odometry.
The mission specifies a world-frame goal $\vect{g}$ and a text query $q$
describing regions to avoid.  No prior geometric map is available.

The planner must output a dynamically feasible local trajectory $\tau_t$.
Following the layered organization in Fig.~\ref{fig:architecture}, mapping,
route selection, and execution compose as
\begin{align}
  \mathcal{G}_t
    &= \Pi_{\mathrm{map}}(\mathcal{G}_{t-1},D_t,\vect{x}_t), \\
  (\mathcal{W}_t,\vect{g}^{\mathrm{front}}_t,
   \vect{g}^{\mathrm{loc}}_t)
    &= \Pi_{\mathrm{route}}(\mathcal{G}_t,I_{0:t},q,\vect{g}), \\
  \tau_t
    &= \Pi_{\mathrm{exec}}(D_t,\vect{o}_t).
\end{align}
Each layer has a deliberately narrow interface.  $\Pi_{\mathrm{map}}$ is the
plug-and-play boundary: it consumes only depth, odometry, and its previous
graph state, and it exposes $\mathcal{G}_t$ without assumptions about the
local policy.  The slower $\Pi_{\mathrm{route}}$ produces the accepted
world-frame witness $\mathcal{W}_t$ together with two goals: a route-level
frontier goal $\vect{g}^{\mathrm{front}}_t$ and a control-level local goal
$\vect{g}^{\mathrm{loc}}_t$ tracked along the witness.  The fast
$\Pi_{\mathrm{exec}}$ is any executor that accepts these standard inputs; we
demonstrate it with the released YOPO policy \cite{yopo}, used without
modification or retraining, which receives one $160\times96$ depth channel,
body-frame motion, and the frontier goal:
\begin{equation}
 \vect{o}_t = [\vect{v}^{B}_t,\vect{a}^{B}_t,
                (\vect{g}^{\mathrm{front}}_t-\vect{p}_t)^{B}] \in \R^9.
\end{equation}
The local policy neither creates new frontiers nor changes the global
homotopy: it chooses a dynamically feasible primitive \emph{inside} the
accepted route.  This division of labor mirrors hierarchical compositions of
topological routes with local control \cite{meng2020,ving}, but transfers
routes as world-frame witness goals rather than as appearance subgoals, and
---unlike those compositions---requires no change to the executor at all.

We use a fixed-height planning layer in the current experiments to isolate
the horizontal branch-selection problem.  Depth sensing, point memory, and
local trajectory execution remain 3-D, and the semantic cost is defined on
3-D witness polylines without requiring the fixed-height assumption.

%%%%%%%%%%%%%%%%%%%%%%%%%%%%%%%%%%%%%%%%%%%%%%%%%%%%%%%%%%%%%%%%%%%%%%%%%%%%%%%%
\section{PLUG-AND-PLAY PERSISTENT MAPPING}
\label{sec:topology}
This layer defines the reusable mapping boundary of ScaleNav.  It requires a
registered depth or point-cloud stream and odometry, and publishes an atomic
graph snapshot containing persistent node identities, verified free-space
bubbles, and edge witnesses.  Route objectives and local-policy features stay
outside the mapper, so the same graph interface can serve a different route
objective or a different executor---the property we mean by
\emph{plug-and-play}.

\subsection{Bounded Geometric Memory}
Depth is back-projected into obstacle points and free-ray evidence.  The
implementation voxel-filters the current frame and retains at most
$N_{\max}$ points; no points from earlier frames are kept in this buffer, and
no occupancy voxels or ESDF are constructed.  This stands in deliberate
contrast to metric memories \cite{octomap,voxblox}: the map's value is its
topology and its verified witnesses, not its volumetric fidelity.

Following prior bubble-graph construction \cite{epic,bubble}, collision-free
bubbles are clustered into graph vertices.  A vertex is
\begin{equation}
 v_i=(\vect{c}_i,\rho_i,z_i,h_i),
\end{equation}
where $\vect{c}_i$ is its center, $\rho_i$ its free-space clearance, $z_i$
its semantic state, and $h_i$ a persistent identity record.  An edge stores a
collision-checked witness polyline
\begin{equation}
 \pi_{ij}=\{\vect{w}^{ij}_0,\ldots,\vect{w}^{ij}_{M_{ij}}\},
 \quad
 L_{ij}=\sum_{m=1}^{M_{ij}}
 \|\vect{w}^{ij}_m-\vect{w}^{ij}_{m-1}\|_2.
\end{equation}
The witness matters twice: it is the executable geometric connection, and it
is the support on which semantic distance is later measured
(Sec.~\ref{sec:semantics}).  The chord $\vect{c}_j-\vect{c}_i$ is
insufficient because a valid witness may bend around an obstacle.

Geometry is updated in a background graph; after local bubble generation,
region clustering, and graph connection, an atomic swap exposes the new graph
to planning.  The point budget and spatial radius bound memory, and
transiently missing bubbles receive a short grace period before deletion.

\subsection{Clearance-Aware Edge Cost}
Let $\bar\rho_{ij}$ be the minimum cached clearance of the two endpoint
bubbles and their witness, and $\rho^\star$ the target clearance.  The soft
clearance term is
\begin{equation}
 C_{\mathrm{clr}}(e_{ij})=
 \lambda_{\mathrm{clr}}L_{ij}
 \left(\frac{\rho^\star}{\rho^\star+\bar\rho_{ij}}\right)^2.
\label{eq:clearance}
\end{equation}
It prefers open routes without declaring a narrow but valid edge invalid;
geometric feasibility remains determined by the bubble graph and its witness
checks, echoing the soft-corridor philosophy of \cite{bubble,gcopter}.

\begin{figure}[!t]
\centering
\safefig[width=\columnwidth]{pics/semantic_risk_field_paper.pdf}
\caption{Real logged ScaleNav decision in Map2.  The left panel is the logged
RGB frame with the deployed PEARL heatmap and $5\times3$ pooling grid; red
rings identify high responses whose synchronized depth has no return.  The
right panel is the first graph update after this evidence changes the route,
with the logged persistent graph, far-field nodes, witness exposure, and
selected witness.  The panels are synchronized by the recorded evidence and
graph timestamps.  Verified witnesses and provisional semantic-frontier
connections are rendered separately; the latter are hypotheses beyond the
measured free-space backbone.}
\label{fig:risk_field}
\end{figure}

%%%%%%%%%%%%%%%%%%%%%%%%%%%%%%%%%%%%%%%%%%%%%%%%%%%%%%%%%%%%%%%%%%%%%%%%%%%%%%%%
\section{SLOW SEMANTIC ROUTE LAYER}
\label{sec:semantics}
\subsection{Semantic Observation Model}
The frozen PEARL model \cite{pearl} maps $(I_t,q)$ to a patch heatmap, which
is max-pooled on a $5\times3$ grid.  A lower-quartile frame baseline is
subtracted and capped at $0.25$ so that weak whole-image response is
suppressed without erasing a uniformly high-risk view.  Every retained patch
is projected by the calibrated pinhole model at a fixed optical-axis depth
$\ell_s$, independent of the measured depth:
\begin{equation}
 \vect{y}_{u,t}=\vect{p}_t+\mat{R}_t
 \left(\vect{t}_{\mathrm{cam}}+\ell_s\widetilde{\vect{d}}_u\right),
\label{eq:projection}
\end{equation}
where $\vect{p}_t$ and $\mat{R}_t$ are the body pose,
$\vect{t}_{\mathrm{cam}}$ the camera translation, and
$\widetilde{\vect{d}}_u=[1,x_u,y_u]^\top$ the unnormalized pinhole ray.
Every endpoint therefore has the same optical-axis depth, while off-axis
endpoints have a larger Euclidean range.  Measured depth is deliberately not
consulted by this projection, so that an empty depth return cannot suppress
semantic evidence.  The difference from RayFronts \cite{rayfronts} is
deliberate: RayFronts casts a semantic \emph{ray} because it seeks
target \emph{direction} for exploration and tolerates unknown range, whereas
we place a semantic \emph{point} at fixed range because we need a
\emph{location} for avoidance and a conservative range estimate is safer than
an optimistic one.  A ray would let a distant hazard cast a long, thin shadow
that the route layer could step across; a point keeps the repulsion compact
and forces the route to go around the observed evidence, not through its
uncertainty.
For raw patch score $s_{u,t}$ and frame baseline $b_t$, the deployed
calibration is
\begin{equation}
 \widetilde{s}_{u,t}=
 \operatorname{clip}_{[0,1]}
 \left(\frac{s_{u,t}-\min(b_t,0.25)}
 {1-\min(b_t,0.25)}\right).
\label{eq:calibration}
\end{equation}
The tuple $(s_i,c_i)$ stores calibrated score and confidence on a persistent
node identity.  Multi-row support, field-of-view margin, and suspected ground
response modulate confidence; timestamp gates reject stale frames.  Online
association uses the latest matched observation rather than temporal score
averaging, while graph replacement restores the stored tuple by persistent
identity.

\subsection{Provisional Semantic Frontiers}
The highest-scoring, spatially separated projected points in each frame
become graph nodes with unknown geometry state, with at most 16 candidates
admitted per frame.  Repeat observations within $r_m$ update an existing
identity.  Older unknown nodes remain stored, but only the current semantic
generation participates in route search; verified nodes retain their
associated semantic state.  Unknown history is pruned behind the vehicle and
capped at 512 nodes.

A current unknown node is attached to at most two of its four nearest
verified backbone candidates.  This straight connection is a provisional
frontier hypothesis, not a measured free-space edge: it may terminate a route
search, but it cannot be used to connect two verified corridors through
unknown space.  Before commitment, the complete candidate witness is checked
against all currently known obstacles; a failed provisional connection is
detached.  When measured geometry later grows a bubble at the node position,
the unknown node is replaced by a verified node that retains its semantic
identity.

A far-field semantic node is therefore neither a hard obstacle nor a
free-space certificate.  Low-risk rays can propose a forward frontier, in the
spirit of semantic frontier selection \cite{vlfm,semexp}, whereas high-risk
rays repel nearby witnesses through the field below.

\subsection{Continuous Risk on Witness Paths}
Endpoint-only node costs react too late: a risky semantic node may sit beyond
a fork while neither endpoint of the fork carries a high score.  We therefore
measure every semantic node against the complete witness polyline:
\begin{equation}
 d_{k,ij}=\min_{\vect{x}\in\pi_{ij}}
             \|\vect{y}_k-\vect{x}\|_2.
\label{eq:anchor_distance}
\end{equation}
The measured endpoint risk is
\begin{equation}
 r^{\mathrm{node}}_{ij}=\frac{1}{2}(s_i c_i+s_j c_j).
\end{equation}
Semantic node $k$ induces the edge risk
\begin{equation}
 r^{\mathrm{ff}}_{k,ij}=s_kc_k
 \exp\left(-\frac{d_{k,ij}^2}{2(R_s/2)^2}\right),
 \qquad d_{k,ij}\le R_s.
\label{eq:field}
\end{equation}
The Gaussian bandwidth is tied to the finite support, $\sigma_s=R_s/2$,
rather than exposing two independently tuned spatial scales.  The combined
risk uses a maximum rather than a sum:
\begin{equation}
 r_{ij}=\max\left(r^{\mathrm{node}}_{ij},
                  \max_k r^{\mathrm{ff}}_{k,ij}\right),
\label{eq:risk}
\end{equation}
which keeps risk independent of the arbitrary density of nearby semantic
nodes.  The semantic edge penalty is the barrier
\begin{equation}
 C_{\mathrm{sem}}(e_{ij})=
 \lambda_{\mathrm{sem}}L_{ij}[-\log(\max(\epsilon,1-r_{ij}))].
\label{eq:semantic_cost}
\end{equation}
It is nearly linear at low risk and grows sharply near unit risk, while the
clamp $\epsilon=10^{-3}$ keeps the planner finite for saturated observations.
The factor $L_{ij}$ makes exposure accumulate with travel distance.  We
emphasize the semantics of this term in the sense of \cite{majumdar}: it is
preference shaping over verified alternatives, not a calibrated collision
probability, and it never vetoes or certifies geometry.

%%%%%%%%%%%%%%%%%%%%%%%%%%%%%%%%%%%%%%%%%%%%%%%%%%%%%%%%%%%%%%%%%%%%%%%%%%%%%%%%
\subsection{Rolling Search and Route Memory}
\subsubsection{Edge and Path Objectives}
For an edge $e_{ij}$, the geometric traversal cost stored by the graph is
$W_{ij}$.  The deployed configuration does not discount remembered edges
during a fresh search; route persistence is enforced instead by the
accepted-route gate below.  The complete candidate edge cost is
\begin{equation}
 C(e_{ij})=W_{ij}+C_{\mathrm{clr}}(e_{ij})+C_{\mathrm{sem}}(e_{ij}).
\label{eq:edge_cost}
\end{equation}
For a graph path $P$, geometric and risk terms are accumulated separately:
\begin{align}
 G(P)&=\sum_{e\in P}W_e,\\
 R(P)&=\sum_{e\in P}
       \left[C_{\mathrm{clr}}(e)+C_{\mathrm{sem}}(e)\right].
\end{align}
Keeping them separate matters: terminal selection may weight geometry and
accumulated risk differently, and neither term receives a hidden
remembered-edge discount.

The slow layer maintains three goal levels: the \emph{mission goal}
$\vect{g}$, a route-level \emph{frontier goal} selected inside the local
graph window of radius $R_g$, and a control-level \emph{local goal} on the
accepted witness.  The mission goal may lie outside the currently connected
graph, so we search all reachable non-viewpoint vertices within radius
$R_g$.  Terminal selection is a single soft loss: only disconnection,
collision, insufficient clearance, or a missing witness is a hard rejection.
For candidate $v$, let $d_{\mathrm{rsv}}(v)$ be the reserve shortage toward
the mission goal, $d_{\mathrm{dir}}(v)$ the misalignment with the mission
direction, $d_{\mathrm{fov}}(v)$ the offset from the field of view, and
$d_{\mathrm{sm}}(v)$ the heading discontinuity at the vehicle.  The
frontier-goal score is
\begin{align}
 J(v)={}&\alpha_gG(P_v)+R(P_v)
       +w_{\mathrm{rsv}}d_{\mathrm{rsv}}(v)+w_{\mathrm{dir}}d_{\mathrm{dir}}(v)
       \nonumber\\
       &+w_{\mathrm{fov}}d_{\mathrm{fov}}(v)
       +w_{\mathrm{sm}}d_{\mathrm{sm}}(v).
\label{eq:terminal}
\end{align}
Direction, field-of-view, and reserve terms guide the search but never
disqualify a feasible terminal; once the mission goal enters the local
window, a bubble search validates the final connection within a short range.
Graph A* uses $f=g+h$ with millisecond-scale connection budgets, and the best
recovered path ends at
\begin{equation}
 v_t^\star=\arg\min_{v\in\mathcal{V}_{\mathrm{cand}}}J(v),
\end{equation}
whose terminal becomes the frontier goal.

\subsubsection{Route Identity Across Graph Updates}
Incremental graph replacement can change node objects even when the physical
route is unchanged---a problem shared by all systems that replan over
rebuilt structures \cite{octomap,voxblox}.  ScaleNav therefore stores the
selected witness in world coordinates:
\begin{equation}
 \mathcal{W}_{t-1}=\{\vect{w}_0,\ldots,\vect{w}_K\}.
\end{equation}
The vehicle is projected onto this polyline and only a forward window of
length $H_r$ is retained:
\begin{equation}
 \mathcal{W}^{+}_{t-1}=
 \operatorname{crop}(\mathcal{W}_{t-1},
                     \operatorname{proj}(\vect{p}_t,\mathcal{W}_{t-1}),H_r).
\end{equation}
An edge in the new graph is remembered if its geometry follows this window
within a lateral tolerance.  The incumbent terminal is restored by persistent
node identity first, with nearest-neighbor remapping as a fallback.
Equation (\ref{eq:edge_cost}) ranks a fresh candidate without a historical
discount; persistence enters only when the route gate compares that candidate
with the remapped incumbent.

A challenger replaces the incumbent only through hysteresis: its risk must
improve by more than a margin, or its total cost must drop below a fixed
ratio of the incumbent cost.  A high-risk incumbent is latched until its
route risk falls back below a release threshold, and prefix-preserving route
extensions are exempt from the switch test.  An increase in semantic risk
requests a new candidate evaluation when it exceeds $\Delta_r$:
\begin{equation}
 \mathcal{R}(\mathcal{W}_{t-1};t)-
 \mathcal{R}(\mathcal{W}_{t-1};t-1)>\Delta_r.
\end{equation}
Risk decreases do not release route memory, and the request is consumed at
the next rolling-search opportunity rather than interrupting a valid
committed route immediately.  This one-sided gate suppresses left--right
oscillation while still allowing new language evidence to affect the next
route decision---the mechanism that turns the slow semantic layer's opinions
into stable flight behavior.

\paragraph{Asynchronous contract and complexity}
The geometry, semantic, and control streams run on separate clocks.  Every
RGB, depth, odometry, graph, and semantic record carries a source timestamp,
and the planner consumes a consistent snapshot rather than joining callbacks
by arrival order.  A graph worker builds topology and witnesses from the
latest bounded depth buffer and publishes them through an atomic replacement,
so a planning tick sees either the previous complete graph or the next one,
never a partially rebuilt set of vertices and edges.

At each tick, semantic state is read together with the graph snapshot and the
remembered world-frame witness.  Timestamp and freshness gates reject stale
state, while a valid incumbent remains eligible during a short graph grace
period.  The route gate can consequently preserve a physical corridor across
replacement without retaining pointers into the old graph.  If the incumbent
fails its geometric witness check, the planner removes the route prior before
searching again; semantic state can rank a replacement but cannot make a
failed witness executable.

Let $V$ and $E$ denote the vertices and candidate edges in the local graph,
$M$ the retained semantic nodes, and $Q$ the samples used to measure each
witness.  A naive semantic edge evaluation costs $O(EQM)$, but the planner
restricts nodes to the edge influence radius, so practical work scales with
the local semantic neighborhood rather than the full history.  The bounded
point buffer, graph window, and semantic-node cap give memory growth
$O(N_{\max}+V+E+M)$ independent of flight duration.  These bounds keep the
slow route layer from competing with the fixed-rate local controller.

%%%%%%%%%%%%%%%%%%%%%%%%%%%%%%%%%%%%%%%%%%%%%%%%%%%%%%%%%%%%%%%%%%%%%%%%%%%%%%%%
\section{EXECUTOR-AGNOSTIC ROUTE TRANSFER}
\label{sec:execution}
The route layer's output is deliberately restricted to what any local
executor already consumes: a route-level frontier goal and a control-level
local goal.  There is no corridor adapter, no additional network input, and
no retraining.  This section explains why two goals suffice, and what the
executor is and is not asked to do.

\subsection{Why a Frontier Goal Alone Is Not Enough}
A single frontier goal communicates where to move but not which homotopy to
preserve: an endpoint-conditioned policy can shortcut away from the accepted
witness---the exact failure mode that motivates topological guidance of local
controllers \cite{ving,meng2020}.  ScaleNav therefore publishes a 15-m local
goal $\vect{g}^{\mathrm{loc}}_t$ that slides along the accepted witness ahead
of the vehicle, in addition to the frontier goal
$\vect{g}^{\mathrm{front}}_t$ that marks the route-level branch decision.
The executor tracks the local goal; when the witness is invalidated or
replaced, the route layer re-projects the vehicle onto the new witness and
the local goal follows, so the executor simply chases a moving target it
already knows how to chase.

\subsection{The Demonstrated Executor}
We use the released YOPO policy \cite{yopo} as $\Pi_{\mathrm{exec}}$.  Its
$160\times96$ depth feature map is decoded on a $3\times5$ grid, giving
$N_{\mathrm{prim}}=15$ candidate directions; each candidate predicts a
nine-dimensional terminal state
$[\vect{p}_1,\vect{v}_1,\vect{a}_1]\in\R^9$ and one scalar score, and a
fixed-duration boundary-value solver converts the state to a smooth
polynomial primitive.  \textbf{ScaleNav changes neither this solver, nor the
candidate set, nor the network weights, nor the training pipeline.}  The only
difference from the baseline condition is the provenance of the goal: the
baseline receives a reactive frontier goal, while ScaleNav supplies a goal
carried by the persistent, language-conditioned witness.

\subsection{What This Buys and What It Costs}
Because the executor is untouched, its published safety behavior and its
primitive feasibility analysis carry over verbatim, and any executor exposing
the same goal interface can be substituted without further work---including
the stronger local layers surveyed in the related work.  The cost of
giving up policy-level conditioning is that fine-grained corridor adherence
is delegated to the route layer: the local goal must be re-issued promptly
when the witness updates, and the witness itself must stay within the
executor's reliable depth envelope.  Both properties are provided by the
route memory and the hysteretic gate of Sec.~\ref{sec:semantics}, and the
closed-loop results of Sec.~\ref{sec:experiments} show that the unmodified
executor, driven this way, retains its perfect safety record while completing
missions faster than its frontier-only counterpart.

%%%%%%%%%%%%%%%%%%%%%%%%%%%%%%%%%%%%%%%%%%%%%%%%%%%%%%%%%%%%%%%%%%%%%%%%%%%%%%%%
\section{EXPERIMENTS}
\label{sec:experiments}
The evaluation mirrors the three-layer design, because each layer makes a
different claim and deserves a different falsification:
\textbf{(Q1)} does the complete stack fly faster and no less safely than a
frontier-only learned policy in closed loop;
\textbf{(Q2)} which components---semantic front end, far-field nodes, route
memory---carry the benefit;
\textbf{(Q3)} does the persistent graph actually preserve routes and reduce
repeated work across legs.
Closed-loop and round-trip trials use UE Map2 through Colosseum/AirSim
\cite{airsim} with synchronized RGB-D and odometry and a 1.6-m planning
layer.  One-way trials fly from $(0,0,1.6)$ to $(0,140,1.6)$~m.

\subsection{Evaluation Protocol}
\label{sec:protocol}
\emph{Statistics.}  Each aggregate row currently contains ten independent
resets.  Where batches are seed matched (the language conditions of
Sec.~\ref{sec:language}), we additionally report paired comparisons; where
they are not (existing baseline and ablation batches), we report sample
statistics without paired significance claims and mark causal conclusions as
pending.  Upgrading all conditions to seed-paired reruns with Wilcoxon
signed-rank tests at $\alpha=0.05$ is scheduled; tables whose cells are
reserved for those reruns show \blank{}, which denotes \emph{pending data},
never a zero.  Failures remain in outcome rates, and continuous flight
metrics are computed over successful trials only.

\emph{Metrics.}  Closed-loop metrics are success, collision, timeout, path
length $L$, flight time $T$, $\bar v=L/T$, maximum speed, and success
weighted by path length (SPL) \cite{anderson2018}.  Language trials add
route-change rate, query compliance, decision distance, switch count, and the
semantic exposure
$E_{\mathrm{sem}}=L^{-1}\int_0^T r(\vect{p}(t))v(t)\,dt$.
Profiling separates point processing, graph construction, A*, the planning
tick, and PEARL latency.

\emph{Baseline fairness.}  EGO-Planner \cite{egoplanner} and SUPER
\cite{super} are optimization stacks designed around their own mapping and
sensing pipelines; porting them to the same monocular RGB-D stream and the
same control stack is itself a configuration problem, and recent benchmarking
shows that such pipelines are highly sensitive to scenario difficulty,
latency, and field-of-view settings \cite{flightbench}.  We therefore report
them as deployed under their recommended configurations, read their failures
in Map2 as configuration- and regime-dependent rather than as absolute
capability limits, and complement the nominal batch with a speed-matched
SUPER batch at 5~m/s.

\subsection{Closed-Loop System Evaluation (Q1)}
\subsubsection{Methods and Conditions}
We compare frontier-only YOPO \cite{yopo}, EGO-Planner \cite{egoplanner},
SUPER \cite{super}, and full ScaleNav.  Language conditions hold geometry and
planner weights fixed while changing only the query; ablations retain the
same primitive generator and state.  PEARL \cite{pearl} runs in all
semantic-enabled conditions.  Table~\ref{tab:params} lists the shared
ScaleNav configuration.

\begin{table}[t]
\caption{Fixed Parameters Used by the Current Implementation}
\label{tab:params}
\centering
\small
\resizebox{\columnwidth}{!}{%
\begin{tabular}{lc@{\quad}lc}
\toprule
Parameter & Value & Parameter & Value \\
\midrule
Point voxel & 0.10 m & Point budget & 100,000 \\
Frame history & current frame only & Graph update & 100 ms \\
Graph rebuild & 100 ms & Search radius & 45 m \\
Local lookahead & 15 m & Goal-connect budget & 20 ms \\
$\lambda_{\rm sem}$ & 2.0 & $R_s$ & 5 m \\
$\lambda_{\rm clr}$ & 2.0 & $\rho^\star$ & 1.2 m \\
Score update & latest match & Edge-memory discount & disabled \\
Optical depth $\ell_s$ & 30 m & Node merge distance & 2.5 m \\
Semantic candidates/frame & 16 & Virtual-node cap & 512 \\
Semantic candidates/edge & 8 & Semantic grid & $5\times3$ \\
Route bubbles $K$ & 12 & Bubble features & center + radius \\
\bottomrule
\end{tabular}
}
\end{table}

\subsubsection{Scenario}
Map2 contains isolated and dense trees, long walls, buildings, and
mixed-scale clusters (Fig.~1(b)).  Its truth export spans
$140.0\times132.0$~m.  In the evaluation corridor $x\in[-45,45]$~m,
$y\in[0,120]$~m, $3.91\%$ of 0.5-m cells intersect truth geometry in the
navigation slab $z\in[0.6,2.6]$~m.  These statistics come from UE truth,
never from any method's partial online map.  Matched language trials
additionally include the unseen query ``avoid power lines.''

\subsubsection{Results}
Figure~\ref{fig:map2_speed_trajectories} shows one recorded run per method on
common truth geometry; incomplete baselines end at their recorded collisions.

\begin{figure}[!ht]
\centering
\safefig[width=0.72\columnwidth]{pics/experiments/map2_0_140_1p6/trajectory_speed_comparison.pdf}
\vspace{-2pt}
\caption{Recorded Map2 trajectories from $(0,0,1.6)$~m to $(0,140,1.6)$~m
over the complete UE Map2 truth map, voxelized at 0.3~m with one-cell display
dilation.  Color denotes executed speed; circles, stars, and crosses mark
starts, completed goals, and incomplete collision endpoints, respectively.}
\label{fig:map2_speed_trajectories}
\end{figure}

\begin{table}[t]
\caption{Ten-Trial Map2 Closed-Loop Comparison Using Conventional Flight
Metrics.  Continuous entries are mean $\pm$ standard deviation over
successes; \blank{} denotes no successful trial in that batch.}
\label{tab:single_run_validation}
\centering
\small
\resizebox{\columnwidth}{!}{%
\begin{tabular}{lccccccc}
\toprule
Method & Success (\%) & Collision (\%) & Timeout (\%) & Path (m) & Time (s) & Avg. speed (m/s) & Max. speed (m/s) \\
\midrule
YOPO            & 100 & 0 & 0 & $158.678\pm0.169$ & $42.210\pm0.128$ & $3.759\pm0.009$ & $5.913\pm0.060$ \\
EGO-Planner     & 0   & 100 & 0 & \blank & \blank & \blank & \blank \\
SUPER & 0 & 80 & 20 & \blank & \blank & \blank & \blank \\
SUPER (5 m/s) & 10 & 60 & 30 & 144.658 & 30.614 & 4.725 & 5.082 \\
\textbf{ScaleNav (ours)} & \textbf{100} & \textbf{0} & \textbf{0} & $157.403\pm5.194$ & $35.495\pm1.752$ & $\mathbf{4.439\pm0.118}$ & $5.221\pm0.165$ \\
\bottomrule
\end{tabular}
}
\end{table}

Table~\ref{tab:single_run_validation} reports ten trials per setting; the
5~m/s SUPER row is a separate batch in which one startup failure was
replaced, and its single success has no dispersion estimate.  ScaleNav
completed all trials with the query ``tree, blocks, wall, line,'' finishing
$18.1\%$ faster in mean flight time than YOPO at comparable path length.  The
EGO-Planner and SUPER outcomes reflect the porting caveats of
Sec.~\ref{sec:protocol}: both stacks were run at their recommended settings
against the shared depth stream, and we do not interpret their Map2 failure
rates as universal.  What the table \emph{does} establish is that, in this
regime, adding the semantic route layer above a learned executor improves
speed without sacrificing the executor's perfect safety record.

\subsection{Language-Conditioned Routing (Q2a)}
\label{sec:language}
The central semantic claim---that language evidence changes routes early and
in the intended direction---requires seed-matched conditions in which only
the query varies.  Ten paired trials per query hold geometry, initial state,
weights, and seed fixed; all cells below are reserved for that replay and
\blank{} marks pending data rather than zero
(Tables~\ref{tab:language_comparison} and~\ref{tab:language_decisions}).
The decision metrics in Table~\ref{tab:language_decisions} are designed as
a \emph{planner-agnostic} evaluation of the semantic route layer, in the
spirit of \cite{rayfronts}'s planner-agnostic search-volume metric: route
change, compliance, decision distance, and semantic exposure are properties
of the route the layer commits to, not of any downstream controller, so they
isolate whether the semantic evidence itself---rather than a lucky local
reaction---produced the avoidance.

\begin{table}[t]
\caption{Language-Condition Comparison, Flight Metrics (Ten Paired Trials,
Seed-Matched; \blank{} = Pending Replay)}
\label{tab:language_comparison}
\centering
\small
\resizebox{\columnwidth}{!}{%
\begin{tabular}{lcccccccc}
\toprule
Query & Succ. (\%) & Coll. (\%) & T/O (\%) & Path (m) & Time (s) & $\bar v$ (m/s) & $v_{\max}$ (m/s) & SPL \\
\midrule
No query & \blank & \blank & \blank & \blank & \blank & \blank & \blank & \blank \\
Avoid buildings & \blank & \blank & \blank & \blank & \blank & \blank & \blank & \blank \\
Avoid trees & \blank & \blank & \blank & \blank & \blank & \blank & \blank & \blank \\
Avoid power lines (unseen) & \blank & \blank & \blank & \blank & \blank & \blank & \blank & \blank \\
\bottomrule
\end{tabular}
}
\end{table}

\begin{table}[t]
\caption{Language-Condition Comparison, Decision Metrics (Ten Paired Trials;
\blank{} = Pending Replay).  Compliance: fraction of trials whose executed
route avoids the queried class by $\ge d_{\mathrm{cmp}}$; decision distance:
range from first high-risk semantic observation to the committed route
change.}
\label{tab:language_decisions}
\centering
\small
\resizebox{\columnwidth}{!}{%
\begin{tabular}{lccccc}
\toprule
Query & Route change (\%) & Compliance (\%) & Decision dist. (m) & Switches & $E_{\mathrm{sem}}$ \\
\midrule
No query & \blank & \blank & \blank & \blank & \blank \\
Avoid buildings & \blank & \blank & \blank & \blank & \blank \\
Avoid trees & \blank & \blank & \blank & \blank & \blank \\
Avoid power lines (unseen) & \blank & \blank & \blank & \blank & \blank \\
\bottomrule
\end{tabular}
}
\end{table}

\subsection{Component Ablation (Q2b)}
Table~\ref{tab:component_ablation} retains the available independent 6~m/s
batches and reserves rows for the pending matched ablations.  These batches
are condition checks, not paired causal estimates; the ``no target''
condition keeps PEARL active with an absent queried object.

\begin{table}[t]
\caption{ScaleNav Variant Batches (Ten Trials Each; \blank{} = Pending
Matched Rerun).  SE: semantic front end; FF: far-field semantic nodes; RM:
world-frame route memory.  The no-target condition retains the semantic front
end.  All variants drive the same unmodified executor.}
\label{tab:component_ablation}
\centering
\small
\resizebox{\columnwidth}{!}{%
\begin{tabular}{@{}l ccc ccccccc@{}}
\toprule
& \multicolumn{3}{c}{Component} & \multicolumn{7}{c}{Closed-Loop Performance} \\
\cmidrule(lr){2-4} \cmidrule(lr){5-11}
Variant & SE & FF & RM & Succ. (\%)$\uparrow$ & Coll. (\%)$\downarrow$ & T/O (\%)$\downarrow$ & Path (m) & Time (s) & $\bar v$ (m/s) & $v_{\max}$ (m/s) \\
\midrule
\textbf{Full ScaleNav} & \checkmark & \checkmark & \checkmark & 100 & 0 & 0 & 157.403$\pm$5.194 & 35.495$\pm$1.752 & 4.439$\pm$0.118 & 5.221$\pm$0.165 \\
\midrule
w/o semantic front end  & $\times$   & $\times$   & \checkmark & 70 & 30 & 0 & 156.510$\pm$3.339 & 33.620$\pm$0.856 & 4.656$\pm$0.071 & 5.445$\pm$0.121 \\
No task-relevant semantics & \checkmark & \checkmark & \checkmark & 80 & 20 & 0 & 154.205$\pm$1.591 & 49.981$\pm$1.247 & 3.087$\pm$0.093 & 4.491$\pm$0.290 \\
w/o far-field nodes     & \checkmark & $\times$   & \checkmark & \blank & \blank & \blank & \blank & \blank & \blank & \blank \\
w/o route memory        & \checkmark & \checkmark & $\times$   & \blank & \blank & \blank & \blank & \blank & \blank & \blank \\
\bottomrule
\end{tabular}
}
\end{table}

The no-semantic batch succeeded in seven trials; earlier batches succeeded in
eight and ten.  This repeat spread precludes a causal semantic claim before
the seed-matched reruns, which we state explicitly rather than paper over.

\subsection{Parameter Sensitivity}
\label{sec:sensitivity}
The discussion of Sec.~\ref{sec:discussion} argues that $\lambda_{\rm sem}$
and $R_s$ have distinct monotone roles; that argument must be backed by
paired evidence.  Table~\ref{tab:sensitivity} reserves a fixed-snapshot
replay over $\lambda_{\rm sem}\in\{0,1,2,4\}$ and $R_s\in\{3,5,7\}$~m with
$\sigma_s=R_s/2$, holding graph snapshots and seeds fixed.  The
$\lambda_{\rm sem}=2$ versus $0$ pair separates semantic preference from
geometric sparsity; the diagonal scans the operating range.

\begin{table}[t]
\caption{Reserved Sensitivity Replay over Fixed Graph Snapshots (Paired
Seeds; \blank{} = Pending).  Path length and exposure are reported as change
relative to the $\lambda_{\rm sem}{=}2$, $R_s{=}5$~m nominal.}
\label{tab:sensitivity}
\centering
\small
\resizebox{\columnwidth}{!}{%
\begin{tabular}{cccccc}
\toprule
$\lambda_{\rm sem}$ & $R_s$ (m) & Success (\%) & $\Delta$Path (\%) & $\Delta E_{\mathrm{sem}}$ (\%) & Switches \\
\midrule
0 & 5 & \blank & \blank & \blank & \blank \\
1 & 5 & \blank & \blank & \blank & \blank \\
\textbf{2} & \textbf{5} & \multicolumn{4}{c}{nominal (reference)} \\
4 & 5 & \blank & \blank & \blank & \blank \\
\midrule
2 & 3 & \blank & \blank & \blank & \blank \\
2 & 7 & \blank & \blank & \blank & \blank \\
\bottomrule
\end{tabular}
}
\end{table}

\subsection{Persistence and Semantic Exposure (Q3)}
\label{sec:persistence}
\emph{Protocol.}  Graph and semantic memory stay alive during a
$(0,0,1.6)\rightarrow(0,140,1.6)\rightarrow(0,0,1.6)$~m round trip.  Goal
events split the legs; 100-Hz odometry gives time and length, while
insertion, rebuild, and route-switch logs measure reuse.  World-frame clouds
are cropped to $z\in[0.3,3.2]$~m and voxelized at 0.5~m.  We measure
occupied-cell density within 3~m of the trajectory, distance to high-risk
semantic nodes near flight height, and the trajectory fraction within 5~m of
them.

\emph{Observed behavior.}  Both available round trips reused the outbound
graph.  Insertions per rebuild fell by $62.7\%$ and $61.5\%$, and switches
fell 18--12 and 15--10.  Rebuild mean nevertheless rose $21.5\%$ and $17.6\%$
as the working set grew; path length rose $1.3\%$ and $4.6\%$.  On the second
return, mean/P90 density fell $8.2\%/28.6\%$, minimum semantic-node distance
rose from 2.56 to 6.93~m, and exposure within 5~m fell from $8.7\%$ to 0.
This demonstrates reuse and a less cluttered, lower-exposure return leg;
opposite viewpoints and correlated geometry, however, prevent a
semantic-causality claim from these runs alone.  The paired replay of
Table~\ref{tab:persistence_replay}---outbound versus return legs over matched
seeds, with and without route memory---is reserved to isolate that effect.

\begin{table}[t]
\caption{Reserved Paired Round-Trip Replay (Outbound vs.\ Return Legs;
\blank{} = Pending).  Density: occupied 0.5-m cells within 3~m of the
trajectory; sem.\ dist.: minimum distance to high-risk semantic nodes.}
\label{tab:persistence_replay}
\centering
\small
\resizebox{\columnwidth}{!}{%
\begin{tabular}{lccccc}
\toprule
Condition & Insertions/rebuild & Switches & Rebuild (ms) & Density (cells) & Sem.\ dist.\ (m) \\
\midrule
Return, with route memory & \blank & \blank & \blank & \blank & \blank \\
Return, w/o route memory & \blank & \blank & \blank & \blank & \blank \\
\bottomrule
\end{tabular}
}
\end{table}

\subsection{Runtime and Memory Profile}
Table~\ref{tab:online_profile} profiles a separate ten-trial batch.  Rates
are per-trial mean $\pm$ standard deviation; event latencies pool trials.
A* is a subset of the planning tick, not an additional serial stage.

\begin{table}[t]
\caption{ScaleNav Online Resource Profile over Ten Map2 Trials}
\label{tab:online_profile}
\centering
\small
\resizebox{\columnwidth}{!}{%
\begin{tabular}{lrrrr}
\toprule
Stage & Samples & Rate (Hz) & Mean (ms) & P95 (ms) \\
\midrule
Point-cloud update & 2,892 & $7.99\pm0.13$ & 3.06 & 5.43 \\
Incremental graph construction & 1,634 & $4.98\pm0.05$ & 60.60 & 94.94 \\
Topological planning tick & 3,259 & $9.95\pm0.06$ & 25.69 & 61.27 \\
Rolling A* portion & 3,259 & $9.95\pm0.06$ & 11.28 & 21.94 \\
PEARL inference & 652 & $1.994\pm0.008$ & 60.08 & 80.00 \\
\bottomrule
\end{tabular}
}
\end{table}

The final graph has $522.0\pm19.0$ nodes, $1965.6\pm98.7$ edges, and
$81.5\pm17.4$ semantic nodes ($22.0\pm4.0$ verified).  PEARL publishes 652
heatmaps and coalesces $75.7\%$ of input frames under its 2-Hz latest-frame
queue; its 2-Hz cadence is the slowest clock in the system and motivates the
asynchronous, snapshot-based contract.  RSS, peak GPU memory, and YOPO policy
latency were not logged and are reserved instrumentation targets.

\subsection{Planned Cross-Map Generalization}
All closed-loop numbers above come from Map2.  To guard against
overfitting the route layer to one layout, Table~\ref{tab:crossmap} reserves
a seed-paired extension to two further UE environments with different
obstacle statistics, evaluated with the frozen configuration of
Table~\ref{tab:params}.

\begin{table}[t]
\caption{Reserved Cross-Map Generalization (Seed-Paired, Frozen
Configuration; \blank{} = Pending)}
\label{tab:crossmap}
\centering
\small
\resizebox{\columnwidth}{!}{%
\begin{tabular}{llcccc}
\toprule
Map & Method & Success (\%) & $\bar v$ (m/s) & Path (m) & SPL \\
\midrule
Map1 & YOPO / ScaleNav & \blank & \blank & \blank & \blank \\
Map2 & YOPO / ScaleNav & \multicolumn{4}{c}{reported in Table~\ref{tab:single_run_validation}} \\
Map3 & YOPO / ScaleNav & \blank & \blank & \blank & \blank \\
\bottomrule
\end{tabular}
}
\end{table}

%%%%%%%%%%%%%%%%%%%%%%%%%%%%%%%%%%%%%%%%%%%%%%%%%%%%%%%%%%%%%%%%%%%%%%%%%%%%%%%%
\section{DISCUSSION AND LIMITATIONS}
\label{sec:discussion}
ScaleNav separates reusable mapping, semantic route choice, and local
execution, and both boundaries are plug-and-play by design.  The mapping
layer's input and output contracts do not depend on the local policy; the
route layer's output is restricted to the executor's native goal interface.
The demonstrated YOPO executor is used entirely unmodified, and any executor
that accepts a frontier goal and a tracked local goal---learned or
optimization-based \cite{egoplanner,gcopter}---can be substituted without
retraining or adapters.  The price of this agnosticism is granularity: the
route layer can steer the executor only through goal placement, so behaviors
that require shaping the primitive distribution itself remain out of reach.

Far-field nodes are optimistic frontier hypotheses, not safety certificates.
Their candidate witnesses are checked against known geometry, while the
safety of unobserved space still depends on subsequent depth and local
avoidance.  This is the same trust asymmetry exploited by beyond-range
frontier methods \cite{rayfronts}, but applied to repulsion rather than
attraction, and it is what allows ScaleNav to act on a power line that depth
has not yet returned without ever claiming that the line's neighborhood is
mapped.

The semantic parameters have distinct and monotone roles.  Increasing
$\lambda_{\rm sem}$ strengthens preference for lower-risk routes without
changing their geometric validity, whereas increasing $R_s$ expands the
influence region of each observation and simultaneously broadens its tied
Gaussian bandwidth.  Excessive values can therefore produce conservative
detours, while small values postpone the branch decision.  The fixed-snapshot
sensitivity protocol of Sec.~\ref{sec:sensitivity} is designed to establish
how wide the stable operating range is; until those replay cells are
populated, robustness to these parameters should not be inferred from the
nominal setting alone.

SUPER \cite{super} and ScaleNav address different layers: SUPER provides
geometric, safety-assured high-speed execution, whereas ScaleNav provides
persistent, language-conditioned route choice.  Their composition---ScaleNav
routes feeding SUPER corridors---is a promising and unevaluated combination.

Finally, the current topological search uses a fixed-height layer; a fully
3-D closed-loop evaluation is future work.  The semantic
interface is model-independent---boxes, points, or masks from detectors such
as Grounding-DINO-class models can be rasterized into the same heatmap
contract---but only PEARL \cite{pearl} is evaluated here, all evaluation is
in simulation through AirSim/Colosseum \cite{airsim}, and the 2-Hz semantic
cadence, while sufficient at 4--5~m/s cruise, bounds the speed at which
semantic foresight remains ahead of the vehicle.

%%%%%%%%%%%%%%%%%%%%%%%%%%%%%%%%%%%%%%%%%%%%%%%%%%%%%%%%%%%%%%%%%%%%%%%%%%%%%%%%
\section{CONCLUSION}
ScaleNav organizes language-guided aerial navigation as three explicit
layers with narrow contracts.  A plug-and-play persistent mapper turns depth
and odometry into reusable, witness-verified free-space topology.  A slow
route layer combines this backbone with provisional far-field semantic
frontiers, continuous witness-level risk, and world-frame route memory
stabilized by a one-sided hysteretic gate.  An unmodified off-the-shelf
executor receives the selected witness as a frontier goal plus a tracked
local goal and retains full control of the local primitive.  The resulting
division is easy
to state and, we argue, the right one for combining foundation-model
perception with flight: mapping provides persistent structure, the route
layer decides where to go, the execution layer decides how to fly, and
semantics may persuade but never certify.  Reserved seed-matched replays will
convert the remaining condition checks into paired causal estimates.

%%%%%%%%%%%%%%%%%%%%%%%%%%%%%%%%%%%%%%%%%%%%%%%%%%%%%%%%%%%%%%%%%%%%%%%%%%%%%%%%
{\footnotesize
\begin{thebibliography}{99}

\bibitem{yopo}
J. Lu, X. Zhang, H. Shen, L. Xu, and B. Tian,
``You Only Plan Once: A Learning-Based One-Stage Planner With Guidance
Learning,'' \emph{IEEE Robot. Autom. Lett.}, vol. 9, no. 7,
pp. 6083--6090, 2024.

\bibitem{rappids}
N. Bucki, J. Lee, and M. W. Mueller,
``Rectangular Pyramid Partitioning Using Integrated Depth Sensors (RAPPIDS):
A Fast Planner for Multicopter Navigation,'' \emph{IEEE Robot. Autom. Lett.},
vol. 5, no. 3, pp. 4626--4633, 2020.

\bibitem{nanomap}
P. R. Florence, J. Carter, J. Ware, and R. Tedrake,
``NanoMap: Fast, Uncertainty-Aware Proximity Queries With Lazy Search of
Local 3-D Data,'' in \emph{Proc. IEEE Int. Conf. Robot. Autom.}, 2018,
pp. 7631--7638.

\bibitem{egoplanner}
X. Zhou, Z. Wang, H. Ye, C. Xu, and F. Gao,
``EGO-Planner: An ESDF-Free Gradient-Based Local Planner for Quadrotors,''
\emph{IEEE Robot. Autom. Lett.}, vol. 6, no. 2, pp. 478--485, 2021.

\bibitem{loquercio}
A. Loquercio, E. Kaufmann, R. Ranftl, M. M\"uller, V. Koltun, and
D. Scaramuzza,
``Learning High-Speed Flight in the Wild,''
\emph{Sci. Robot.}, vol. 6, no. 59, 2021.

\bibitem{kaufmann2023}
E. Kaufmann, L. Bauersfeld, A. Loquercio, M. M\"uller, V. Koltun, and
D. Scaramuzza,
``Champion-Level Drone Racing Using Deep Reinforcement Learning,''
\emph{Nature}, vol. 620, no. 7976, pp. 982--987, 2023.

\bibitem{bubble}
Y. Ren, F. Zhu, W. Liu, Z. Wang, Y. Lin, F. Gao, and F. Zhang,
``Bubble Planner: Planning High-Speed Smooth Quadrotor Trajectories Using
Receding Corridors,'' in \emph{Proc. IEEE/RSJ Int. Conf. Intell. Robots
Syst.}, 2022, pp. 6332--6339.

\bibitem{epic}
S. Geng, Z. Ning, F. Zhang, and B. Zhou,
``EPIC: A Lightweight LiDAR-Based AAV Exploration Framework for Large-Scale
Scenarios,'' \emph{IEEE Robot. Autom. Lett.}, vol. 10, no. 5,
pp. 5090--5097, 2025.

\bibitem{sfc}
S. Liu, M. Watterson, K. Mohta, K. Sun, S. Bhattacharya, C. J. Taylor, and
V. Kumar,
``Planning Dynamically Feasible Trajectories for Quadrotors Using Safe Flight
Corridors in 3-D Complex Environments,'' \emph{IEEE Robot. Autom. Lett.},
vol. 2, no. 3, pp. 1688--1695, 2017.

\bibitem{gcopter}
Z. Wang, X. Zhou, C. Xu, and F. Gao,
``Geometrically Constrained Trajectory Optimization for Multicopters,''
\emph{IEEE Trans. Robot.}, vol. 38, no. 5, pp. 3259--3278, 2022.

\bibitem{fastplanner}
X. Zhou, J. Gao, C. Wang, Z. Wang, H. Ye, and F. Gao,
``Fast-Planner: A Robust and Efficient Trajectory Planner for Autonomous
Quadrotor Flight,'' in \emph{Proc. IEEE/RSJ Int. Conf. Intell. Robots Syst.},
2019, pp. 8455--8461.

\bibitem{faster}
J. Tordesillas, B. T. Lopez, J. P. How, and L. Carlone,
``FASTER: A Robust Fast-Planner for Safe Autonomous Navigation in Unknown
Environments,'' \emph{IEEE Trans. Robot.}, vol. 38, no. 2, pp. 1112--1131,
2022.

\bibitem{super}
Y. Ren, F. Zhu, G. Lu, Y. Cai, L. Yin, F. Kong, J. Lin, N. Chen, and
F. Zhang,
``Safety-Assured High-Speed Navigation for MAVs,''
\emph{Sci. Robot.}, vol. 10, no. 98, p. eado6187, 2025.

\bibitem{yamauchi}
B. Yamauchi,
``A Frontier-Based Approach for Autonomous Exploration,'' in \emph{Proc.
IEEE Int. Symp. Comput. Intell. Robot. Autom.}, 1997, pp. 146--151.

\bibitem{fuel}
B. Zhou, Y. Zhang, X. Chen, and S. Shen,
``FUEL: Fast UAV Exploration Using Incremental Frontier Structure and
Hierarchical Planning,'' \emph{IEEE Robot. Autom. Lett.}, vol. 6, no. 2,
pp. 779--786, 2021.

\bibitem{tare}
C. Cao, H. Zhu, H. Choset, and J. Zhang,
``TARE: A Hierarchical Framework for Efficiently Exploring Complex 3D
Environments,'' in \emph{Proc. Robot.: Sci. Syst.}, 2021.

\bibitem{octomap}
K. M. Wurm, A. Hornung, M. Bennewitz, C. Stachniss, and W. Burgard,
``OctoMap: An Efficient Probabilistic 3D Mapping Framework Based on
Octrees,'' \emph{Auton. Robots}, vol. 34, pp. 189--206, 2013.

\bibitem{voxblox}
H. Oleynikova, Z. Taylor, M. Fehr, R. Siegwart, and J. Nieto,
``Voxblox: Incremental 3D Euclidean Signed Distance Fields for On-Board MAV
Planning,'' \emph{IEEE/RSJ Int. Conf. Intell. Robots Syst. (IROS)}, 2017.

\bibitem{kimera}
A. Rosinol, A. Violette, M. Abate, N. Hughes, Y. Chang, J. Shi, A. Gupta, and
L. Carlone,
``Kimera: From SLAM to Spatial Perception With 3D Dynamic Scene Graphs,''
\emph{Int. J. Robot. Res.}, vol. 40, no. 12--14, pp. 1510--1546, 2021.

\bibitem{hydra}
N. Hughes, Y. Chang, and L. Carlone,
``Hydra: A Real-Time Spatial Perception System for 3D Scene Graph
Construction and Optimization,'' in \emph{Proc. Robot.: Sci. Syst.}, 2022.

\bibitem{sptm}
N. Savinov, A. Dosovitskiy, and V. Koltun,
``Semi-Parametric Topological Memory for Navigation,'' in \emph{Proc. Int.
Conf. Learn. Represent.}, 2018.

\bibitem{ving}
D. Shah, B. Eysenbach, G. Kahn, N. Rhinehart, and S. Levine,
``ViNG: Learning Open-World Navigation With Visual Goals,'' in \emph{Proc.
IEEE Int. Conf. Robot. Autom.}, 2021, pp. 13215--13222.

\bibitem{viking}
D. Shah and S. Levine,
``ViKiNG: Vision-Based Kilometer-Scale Navigation With Geographic Hints,''
in \emph{Proc. Robot.: Sci. Syst.}, 2022.

\bibitem{vint}
D. Shah, A. Sridhar, N. Dashora, K. Stachowicz, K. Black, N. Hirose, and
S. Levine,
``ViNT: A Foundation Model for Visual Navigation,'' in \emph{Proc. Conf.
Robot Learn.}, 2023.

\bibitem{nomad}
A. Sridhar, D. Shah, C. Glossop, and S. Levine,
``NoMaD: Goal Masked Diffusion Policies for Navigation and Exploration,''
in \emph{Proc. IEEE Int. Conf. Robot. Autom.}, 2024, pp. 63--70.

\bibitem{meng2020}
X. Meng, N. Ratliff, Y. Xiang, and D. Fox,
``Scaling Local Control to Large-Scale Topological Navigation,'' in
\emph{Proc. IEEE Int. Conf. Robot. Autom.}, 2020, pp. 672--678.

\bibitem{clip}
A. Radford et al.,
``Learning Transferable Visual Models From Natural Language Supervision,''
in \emph{Proc. Int. Conf. Mach. Learn.}, 2021, pp. 8748--8763.

\bibitem{lseg}
B. Li, K. Q. Weinberger, S. Belongie, V. Koltun, and R. Ranftl,
``Language-Driven Semantic Segmentation,'' in \emph{Proc. Int. Conf. Learn.
Represent.}, 2022.

\bibitem{openseg}
G. Ghiasi, X. Gu, Y. Cui, and T.-Y. Lin,
``Scaling Open-Vocabulary Image Segmentation With Image-Level Labels,'' in
\emph{Proc. Eur. Conf. Comput. Vis.}, 2022.

\bibitem{pearl}
G. Pei, X. Jiang, X. Cai, T. Chen, Y. Yao, and B. Jeon,
``PEARL: Geometry Aligns Semantics for Training-Free Open-Vocabulary Semantic
Segmentation,'' \emph{arXiv preprint arXiv:2603.21528}, 2026.

\bibitem{lerf}
J. Kerr, C. M. Kim, K. Goldberg, A. Kanazawa, and M. Tancik,
``LERF: Language Embedded Radiance Fields,'' in \emph{Proc. IEEE/CVF Int.
Conf. Comput. Vis.}, 2023.

\bibitem{clipfields}
N. M. M. Shafiullah, C. Paxton, L. Pinto, S. Chintala, and A. Szlam,
``CLIP-Fields: Weakly Supervised Semantic Fields for Robotic Memory,'' in
\emph{Proc. Robot.: Sci. Syst.}, 2023.

\bibitem{conceptfusion}
K. M. Jatavallabhula et al.,
``ConceptFusion: Open-Set Multimodal 3D Mapping,'' in \emph{Proc. Robot.:
Sci. Syst.}, 2023.

\bibitem{openscene}
S. Peng, K. Genova, C. Jiang, A. Tagliasacchi, M. Pollefeys, and
T. Funkhouser,
``OpenScene: 3D Scene Understanding With Open Vocabularies,'' in
\emph{Proc. IEEE/CVF Conf. Comput. Vis. Pattern Recognit.}, 2023,
pp. 815--824.

\bibitem{conceptgraphs}
Q. Gu et al.,
``ConceptGraphs: Open-Vocabulary 3D Scene Graphs for Perception and
Planning,'' in \emph{Proc. IEEE Int. Conf. Robot. Autom.}, 2024,
pp. 5021--5028.

\bibitem{vlmaps}
C. Huang, O. Mees, A. Zeng, and W. Burgard,
``Visual Language Maps for Robot Navigation,'' in \emph{Proc. IEEE Int.
Conf. Robot. Autom.}, 2023, pp. 10608--10615.

\bibitem{hovsg}
A. Werby, C. Huang, M. B\"uchner, A. Valada, and W. Burgard,
``Hierarchical Open-Vocabulary 3D Scene Graphs for Language-Grounded Robot
Navigation,'' in \emph{Proc. Robot.: Sci. Syst.}, 2024.

\bibitem{blukis}
V. Blukis, C. Paxton, D. Fox, A. Garg, and Y. Artzi,
``A Persistent Spatial Semantic Representation for High-Level Natural
Language Instruction Execution,'' in \emph{Proc. Conf. Robot Learn.}, 2021,
pp. 706--717.

\bibitem{semexp}
D. S. Chaplot, D. Gandhi, A. Gupta, and R. Salakhutdinov,
``Object Goal Navigation Using Goal-Oriented Semantic Exploration,'' in
\emph{Proc. Adv. Neural Inf. Process. Syst.}, 2020, pp. 4247--4258.

\bibitem{vlfm}
N. Yokoyama, S. Ha, D. Batra, J. Wang, and B. Bucher,
``VLFM: Vision-Language Frontier Maps for Zero-Shot Semantic Navigation,''
in \emph{Proc. IEEE Int. Conf. Robot. Autom.}, 2024, pp. 42--48.

\bibitem{lmnav}
D. Shah, B. Osi\'nski, B. Ichter, and S. Levine,
``LM-Nav: Robotic Navigation With Large Pre-Trained Models of Language,
Vision, and Action,'' in \emph{Proc. Conf. Robot Learn.}, 2022,
pp. 492--504.

\bibitem{rayfronts}
O. Alama, A. Bhattacharya, H. He, S. Kim, Y. Qiu, W. Wang, C. Ho, N. Keetha,
and S. Scherer,
``RayFronts: Open-Set Semantic Ray Frontiers for Online Scene Understanding
and Exploration,'' in \emph{Proc. IEEE/RSJ Int. Conf. Intell. Robots Syst.},
2025, pp. 5930--5937.

\bibitem{wildos}
H. Shah, E. Tevere, D. Atha, M. Kaufmann, S. Khattak, M. Patel, M. Hutter,
J. Frey, and P. Spieler,
``WildOS: Open-Vocabulary Object Search in the Wild,''
\emph{arXiv preprint arXiv:2602.19308}, 2026.

\bibitem{aerialvln}
S. Liu, H. Zhang, Y. Qi, P. Wang, Y. Zhang, and Q. Wu,
``AerialVLN: Vision-and-Language Navigation for UAVs,'' in \emph{Proc.
IEEE/CVF Int. Conf. Comput. Vis.}, 2023, pp. 15338--15348.

\bibitem{bartolomei}
L. Bartolomei, L. Teixeira, and M. Chli,
``Semantic-Aware Active Perception for UAVs Using Deep Reinforcement
Learning,'' in \emph{Proc. IEEE/RSJ Int. Conf. Intell. Robots Syst.}, 2021,
pp. 3101--3108.

\bibitem{blackmore}
L. Blackmore, M. Ono, and B. C. Williams,
``Chance-Constrained Optimal Path Planning With Obstacles,'' \emph{IEEE
Trans. Robot.}, vol. 27, no. 6, pp. 1080--1094, 2011.

\bibitem{majumdar}
A. Majumdar and M. Pavone,
``How Should a Robot Assess Risk? Towards an Axiomatic Theory of Risk in
Robotics,'' in \emph{Robotics Research: The 18th International Symposium
(ISRR)}, Springer, 2020, pp. 75--84.

\bibitem{rast}
G. Chen, S. Wu, M. Shi, W. Dong, H. Zhu, and J. Alonso-Mora,
``RAST: Risk-Aware Spatio-Temporal Safety Corridors for MAV Navigation in
Dynamic Uncertain Environments,'' \emph{IEEE Robot. Autom. Lett.}, vol. 8,
no. 2, pp. 808--815, 2023.

\bibitem{anderson2018}
P. Anderson, A. Chang, D. S. Chaplot, A. Dosovitskiy, S. Gupta, V. Koltun,
J. Kosecka, J. Malik, R. Mottaghi, M. Savva, and A. R. Zamir,
``On Evaluation of Embodied Navigation Agents,'' \emph{arXiv preprint
arXiv:1807.06757}, 2018.

\bibitem{flightbench}
S.-A. Yu, C. Yu, F. Gao, Y. Wu, and Y. Wang,
``FlightBench: Benchmarking Learning-Based Methods for Ego-Vision-Based
Quadrotors Navigation,'' in \emph{Proc. IEEE Int. Conf. Robot. Autom.},
2025.

\bibitem{airsim}
S. Shah, D. Dey, C. Lovett, and A. Kapoor,
``AirSim: High-Fidelity Visual and Physical Simulation for Autonomous
Vehicles,'' in \emph{Field and Service Robotics}, 2018, pp. 621--635.

\end{thebibliography}
}

\end{document}
